\chapter{Comparative Statics}\label{ch:b1:comparative-statics}

A government raises the tax on a litre of fuel by ten cents. Before anyone
fills a tank at the new price, a ministry wants to know how much the pump
price will rise, how much less fuel will be sold, and how much the tax will
raise. A model gives the \omterm{def:b1:models-and-data:equilibrium}{equilibrium} for the old tax; the question is about
how the \omterm{def:b1:models-and-data:equilibrium}{equilibrium} moves when the tax changes. That question, asked of
every model in this book, is \omterm{def:b1:comparative-statics:comparative-statics}{comparative statics}. This chapter gives its
three tools: the derivative of a solution found from an equation it solves,
the elasticity as a derivative of logarithms, and the envelope theorem,
which tells how the best attainable value moves without solving the
problem again.

\begin{openingchart}{Energy use per person against GDP per person, 183 countries,
2019, both on logarithmic axes, with the least-squares line through the
logarithms. Its slope, 1.25, is an elasticity read off a chart.}{U.S. Energy
Information Administration, energy consumption per person (million British
thermal units); World Bank, World Development Indicators, GDP per person at
purchasing-power parity, constant 2021 international dollars}
% figdata: energy_income_stats (slope 1.25, n 183)
\begin{tikzpicture}
\begin{loglogaxis}[omscatter, log ticks with fixed point, xtickten={-3,...,9}, ytickten={-3,...,9},
  xlabel={GDP per person (dollars, logarithmic axis)},
  ylabel={Energy use per person (million Btu)}, xmin=500, xmax=200000,
  ymin=0.1, ymax=1000, height=0.55\linewidth]
  \addplot table[x=gdppc, y=energy] {figdata/bachelor-1/03-comparative-statics/out/energy_income_2019.csv};
  \addplot[omThm, thick, no marks, sharp plot] table[x=gdppc, y=fit]
    {figdata/bachelor-1/03-comparative-statics/out/energy_fit_2019.csv};
\end{loglogaxis}
\end{tikzpicture}
\end{openingchart}

On logarithmic axes, the cloud of 183 countries lies close to a straight
line: the correlation of the two logarithms is $0.94$. % figdata: energy_income_stats
The slope of the line is 1.25. % figdata: energy_income_stats
\Cref{sec:b1:comparative-statics:elasticity} shows that the slope of a
straight line on such a chart is an elasticity: across countries, one per
cent more income per person goes with 1.25 per cent more energy use per
person. India, at \usd{7702} a head, used 21.85 million British thermal
units a person; the United States, nine times richer at \usd{69121}, used
292.7, thirteen times as much. % figdata: energy_income_2019
Whether this number says anything about how one country's energy use will
move when its income moves is a question of what the chart can identify,
and the chapter returns to it in a measure box. The chapter's main concern
is the opposite kind of number: the response of a model's solution to a
change made on purpose, a tax.

\begin{recall}
The price elasticity of demand is the percentage change in quantity
demanded divided by the percentage change in price; demand is elastic when
it is below $-1$. Arc and point elasticities differ in the base of the
percentages. A tax on sellers shifts supply up by the tax; buyers bear the
larger share of it when demand is less elastic than supply.
\end{recall}

\section{How a solution moves}\label{sec:b1:comparative-statics:moves}

A model with one \omterm{def:b1:models-and-data:variables}{endogenous variable} $x$ and one \omterm{def:b1:models-and-data:variables}{exogenous variable} or
parameter $\theta$ can be written as one equation $F(x, \theta) = 0$, whose
solution is $x^*(\theta)$ (\cref{def:b1:models-and-data:variables}).

\begin{definition}[Comparative statics]\label{def:b1:comparative-statics:comparative-statics}
\index{comparative statics}
The \emph{comparative statics} of a model are the derivatives of its
solution with respect to its \omterm{def:b1:models-and-data:variables}{exogenous variables} and parameters,
$\dd x^*/\dd\theta$: how the \omterm{def:b1:models-and-data:equilibrium}{equilibrium} or the optimum moves when one given
variable changes and the others are held fixed.
\end{definition}

\begin{proposition}[Implicit derivative]\label{prop:b1:comparative-statics:implicit}
\index{implicit derivative}\index{comparative statics!implicit derivative}
Let $F$ have continuous partial derivatives, let $x^*(\theta)$ be a
differentiable solution of $F(x, \theta) = 0$, and let
$F_x(x^*, \theta) \neq 0$. Then
\[
\frac{\dd x^*}{\dd\theta} = -\frac{F_\theta(x^*, \theta)}{F_x(x^*, \theta)} .
\]
\end{proposition}

\begin{proof}
$F(x^*(\theta), \theta) = 0$ for every $\theta$. Differentiate with respect to
$\theta$ by the chain rule: $F_x \,\dd x^*/\dd\theta + F_\theta = 0$. Divide by
$F_x \neq 0$.
\end{proof}

That $x^*$ is differentiable when $F_x \neq 0$ is again the implicit
function theorem, admitted. The condition $F_x \neq 0$ matters: where it
fails, a small change in $\theta$ can make the solution jump or vanish. The
proposition is \cref{prop:b1:constrained-optimisation:level-slope} in
another dress: there $\theta$ was $x$ and $x^*$ was $y$.

Apply it to a market with a per-unit tax $\tau$ collected from sellers.
Buyers pay $p$, sellers keep $p - \tau$, demand is $D(p)$ and supply
$S(p - \tau)$. The \omterm{def:b1:models-and-data:equilibrium}{equilibrium} price solves $F(p, \tau) = D(p) - S(p - \tau) =
0$, with $F_p = D'(p) - S'(p - \tau)$ and $F_\tau = S'(p - \tau)$.

\begin{proposition}[Incidence of a per-unit tax]\label{prop:b1:comparative-statics:incidence}
\index{tax incidence}\index{pass-through}
If $D' < 0$ and $S' > 0$, the price paid by buyers, the price kept by
sellers and the quantity move with the tax as
\[
\frac{\dd p}{\dd\tau} = \frac{S'}{S' - D'} \in (0, 1), \qquad
\frac{\dd (p - \tau)}{\dd\tau} = \frac{D'}{S' - D'} \in (-1, 0), \qquad
\frac{\dd q}{\dd\tau} = \frac{D' S'}{S' - D'} < 0 .
\]
\end{proposition}

\begin{proof}
By \cref{prop:b1:comparative-statics:implicit},
$\dd p/\dd\tau = -S'/(D' - S') = S'/(S' - D')$; numerator and denominator
are positive and the numerator is smaller. Subtracting 1 gives the second
expression. The quantity is $q = D(p)$, so $\dd q/\dd\tau = D' \,\dd p/\dd\tau$.
\end{proof}

The fraction $\dd p/\dd\tau$ is the \emph{pass-through} of the tax: the part
of each unit of tax that shows up in the price buyers pay. With
$D(p) = 100 - 2p$ and $S(p - \tau) = 3(p - \tau)$, $D' = -2$ and $S' = 3$, so
the pass-through is $3/5 = 0.6$ and the quantity falls by $1.2$ units per unit
of tax. Here the model can be solved outright, $p = 20 + 0.6\tau$ and
$q = 60 - 1.2\tau$, and the derivatives agree; the proposition gives them
when it cannot.

\begin{omfigure}
\begin{tikzpicture}
\begin{axis}[omrecord, grid=none, name=left, xlabel={Quantity}, ylabel={Price},
  xmin=0, xmax=80, ymin=0, ymax=50, width=0.48\linewidth, height=0.42\linewidth,
  title={\small More elastic demand}, samples=2]
  \addplot[demand, domain=0:80] {50 - x/2};
  \addplot[supply, domain=0:80] {x/3};
  \addplot[supply, shifted, domain=0:80] {x/3 + 5};
  \addplot[only marks, mark=*, black] coordinates {(60,20) (54,23)};
  \addplot[guide] coordinates {(0,23) (54,23) (54,0)};
  \addplot[guide] coordinates {(0,20) (60,20) (60,0)};
\end{axis}
\begin{axis}[omrecord, grid=none, at={(left.east)}, anchor=west, xshift=12mm,
  xlabel={Quantity}, xmin=0, xmax=80, ymin=0, ymax=50, width=0.48\linewidth,
  height=0.42\linewidth, title={\small Less elastic demand}, samples=2]
  \addplot[demand, domain=30:70] {140 - 2*x};
  \addplot[supply, domain=0:80] {x/3};
  \addplot[supply, shifted, domain=0:80] {x/3 + 5};
  \addplot[only marks, mark=*, black] coordinates {(60,20) (57.857,24.286)};
  \addplot[guide] coordinates {(0,24.286) (57.857,24.286)};
  \addplot[guide] coordinates {(0,20) (60,20)};
\end{axis}
\end{tikzpicture}

{\small A tax of 5 on sellers shifts supply $q = 3(p - \tau)$ up by 5 (dashed),
from the same starting point $(60, 20)$. Left, demand $q = 100 - 2p$: the price
rises by 3 and the quantity falls by 6. Right, demand $q = 70 - 0.5p$: the
price rises by 4.29 and the quantity falls by 2.14.}
\end{omfigure}

\section{Elasticities as logarithmic derivatives}\label{sec:b1:comparative-statics:elasticity}

The two panels start from the same point and differ only in the slope of
demand. Slopes depend on units; a model that must be compared across goods
and countries needs a measure that does not. The school volumes used the
ratio of percentage changes. With calculus it becomes a derivative.

\begin{definition}[Point elasticity]\label{def:b1:comparative-statics:point-elasticity}
\index{elasticity!point}\index{point elasticity}
The \emph{elasticity} of a positive differentiable function $y(x)$, with
$x > 0$, is
\[
\varepsilon_{y,x} = \frac{\dd\ln y}{\dd\ln x} = \frac{x}{y}\,\frac{\dd y}{\dd x}.
\]
It is the percentage change in $y$ per percentage change in $x$, in the
limit of small changes, and it has no unit.
\end{definition}

The two expressions agree by the chain rule, since $\dd\ln y = \dd y/y$ and
$\dd\ln x = \dd x/x$. In the left panel at $(60, 20)$, the demand elasticity
is $-2 \times 20/60 = -0.67$; in the right panel, $-0.5 \times 20/60 = -0.17$.
Supply has the elasticity $3 \times 20/60 = 1$ in both.

\begin{proposition}[Pass-through in elasticities]\label{prop:b1:comparative-statics:pass-through}
\index{pass-through!in elasticities}
At a point where $\tau = 0$, with $\varepsilon_d < 0$ and $\varepsilon_s > 0$ the
elasticities of demand and supply there,
\[
\frac{\dd p}{\dd\tau} = \frac{\varepsilon_s}{\varepsilon_s - \varepsilon_d} .
\]
\end{proposition}

\begin{proof}
Multiply the numerator and the denominator of $S'/(S' - D')$ in
\cref{prop:b1:comparative-statics:incidence} by $p/q$: $S'p/q = \varepsilon_s$
and $D'p/q = \varepsilon_d$, since at $\tau = 0$ both curves pass through $(p, q)$.
\end{proof}

Left panel: $1/(1 + 0.67) = 0.6$; right panel: $1/(1 + 0.17) = 0.86$. The less
elastic side of the market bears more of the tax. The proposition uses only
two numbers that can be estimated, which is why the incidence of a tax is
discussed in elasticities.

\begin{definition}[Constant-elasticity function]\label{def:b1:comparative-statics:constant-elasticity}
\index{constant-elasticity function}
A \emph{constant-elasticity function} is $y = A x^{\varepsilon}$ with $A > 0$:
its elasticity is $\varepsilon$ at every $x$, and its graph on logarithmic axes
is the straight line $\ln y = \ln A + \varepsilon \ln x$.
\end{definition}

\begin{proposition}[Rules for elasticities]\label{prop:b1:comparative-statics:rules}
\index{elasticity!rules}
For positive differentiable functions,
$\varepsilon_{yz, x} = \varepsilon_{y,x} + \varepsilon_{z,x}$,
$\varepsilon_{y/z, x} = \varepsilon_{y,x} - \varepsilon_{z,x}$, and if $y$
depends on $x$ through $z$, $\varepsilon_{y,x} = \varepsilon_{y,z}\,\varepsilon_{z,x}$.
\end{proposition}

\begin{proof}
Take logarithms: $\ln(yz) = \ln y + \ln z$ and $\ln(y/z) = \ln y - \ln z$, and
differentiate with respect to $\ln x$. For the third, the chain rule gives
$\dd\ln y/\dd\ln x = (\dd\ln y/\dd\ln z)(\dd\ln z/\dd\ln x)$.
\end{proof}

Spending on a good is $pq$, so its elasticity with respect to the price is
$1 + \varepsilon_d$: spending rises with the price when demand is inelastic,
$\varepsilon_d > -1$, and falls when it is elastic. That is the school rule,
now a one-line consequence of the logarithm.

\begin{omfigure}
\begin{tikzpicture}
\begin{axis}[omrecord, grid=none, name=lin, xlabel={Price $p$}, ylabel={Quantity $q$},
  xmin=0, xmax=25, ymin=0, ymax=60, width=0.46\linewidth, height=0.4\linewidth,
  title={\small Ordinary axes}, samples=80]
  \addplot[demand, domain=0.7:25] {50/sqrt(x)};
\end{axis}
\begin{loglogaxis}[omrecord, grid=none, at={(lin.east)}, anchor=west, xshift=12mm,
  log ticks with fixed point, xtickten={-3,...,9}, ytickten={-3,...,9},
  xlabel={Price $p$}, xmin=0.5, xmax=100, ymin=1, ymax=100,
  width=0.46\linewidth, height=0.4\linewidth, title={\small Logarithmic axes}, samples=2]
  \addplot[demand, domain=0.5:100] {50/sqrt(x)};
\end{loglogaxis}
\end{tikzpicture}

{\small The constant-elasticity demand $q = 50p^{-0.5}$: a curve on ordinary
axes, a straight line of slope $-0.5$ on logarithmic axes.}
\end{omfigure}

\begin{measure}[An elasticity from a logarithmic chart]
The slope of a straight line through a cloud on logarithmic axes is an
elasticity of the line, not necessarily of anything that causes anything.
The opening chart's 1.25 compares countries that differ in climate, in the
size of their territory, in the share of industry in output and in the price
of energy. % figdata: energy_income_stats
It is not the response of one country's energy use to a rise in its own
income, still less the response of one household's. It is a measure of how
energy and income go together across the world in one year, and it can
change: on the 152 countries present in both years, the slope was 1.35 in
1990 and 1.25 in 2019. % figdata: energy_income_common
To turn a slope into an effect, one needs what
\cref{sec:b1:models-and-data:identification} required: a reason why the
variation on the horizontal axis moves only the curve one wants to trace.
\end{measure}

\begin{omchart}{Energy use per person against GDP per person, 1990 and 2019,
logarithmic axes, with the least-squares line of each year.}{U.S. Energy
Information Administration; World Bank, World Development Indicators; as in
the opening chart}
\begin{tikzpicture}
\begin{loglogaxis}[omscatter, log ticks with fixed point, xtickten={-3,...,9}, ytickten={-3,...,9},
  xlabel={GDP per person (dollars, logarithmic axis)},
  ylabel={Energy use per person (million Btu)}, xmin=300, xmax=200000,
  ymin=0.1, ymax=1000, height=0.5\linewidth, legend pos=north west]
  \addplot table[x=gdppc, y=energy] {figdata/bachelor-1/03-comparative-statics/out/energy_income_1990.csv};
  \addplot table[x=gdppc, y=energy] {figdata/bachelor-1/03-comparative-statics/out/energy_income_2019.csv};
  \addplot[omDef, thick, no marks, sharp plot, forget plot] table[x=gdppc, y=fit]
    {figdata/bachelor-1/03-comparative-statics/out/energy_fit_1990.csv};
  \addplot[omThm, thick, no marks, sharp plot, forget plot] table[x=gdppc, y=fit]
    {figdata/bachelor-1/03-comparative-statics/out/energy_fit_2019.csv};
  \legend{1990, 2019}
\end{loglogaxis}
\end{tikzpicture}
\end{omchart}

\section{The envelope theorem}\label{sec:b1:comparative-statics:envelope}

\omterm{def:b1:comparative-statics:comparative-statics}{Comparative statics} of an optimum can be asked of two objects: the choice,
$x^*(\theta)$, and the best attainable value, $V(\theta)$. The second is
easier than it looks.

\begin{theorem}[Envelope theorem]\label{thm:b1:comparative-statics:envelope}
\index{envelope theorem}
\begin{enumerate}
  \item Let $V(\theta) = \max_x f(x, \theta)$, attained at an interior
    $x^*(\theta)$ that is differentiable in $\theta$, with $f$ smooth. Then
    $V'(\theta) = f_\theta\bigl(x^*(\theta), \theta\bigr)$.
  \item Let $V(\theta) = \max f(x, y, \theta)$ subject to
    $g(x, y, \theta) = 0$, with solution $(x^*, y^*)$ and multiplier
    $\lambda^*$ satisfying the \omterm{thm:b1:constrained-optimisation:lagrange}{Lagrange conditions} and differentiable in
    $\theta$. Then
    \[
    V'(\theta) = \pd{\mathcal L}{\theta}\Big|_{(x^*, y^*, \lambda^*)}
    = f_\theta - \lambda^* g_\theta .
    \]
\end{enumerate}
\end{theorem}

\begin{proof}
(1) By the chain rule $V'(\theta) = f_x \,\dd x^*/\dd\theta + f_\theta$, and
$f_x = 0$ at an interior maximum. (2) The proof is in
\cref{exo:b1:comparative-statics:envelope-proof}.
\end{proof}

The theorem says that when a parameter moves, the optimiser's own
adjustment of $x$ has no first-order effect on the value: at the optimum,
small changes in $x$ are worth nothing. Only the direct effect of $\theta$
counts. \Cref{prop:b1:constrained-optimisation:shadow} is the case
$g = g_0(x, y) - b$, $\theta = b$: $\partial\mathcal L/\partial b = \lambda^*$.

The name comes from a picture. A firm whose cost of producing $q$ with a
plant of size $K$ is $C(q, K) = q^2/K + K$ (invented) has, for each plant,
a short-run cost curve. Choosing the best plant for each $q$ gives the
long-run cost $C^{\mathrm{LR}}(q) = \min_K C(q, K)$: the condition
$-q^2/K^2 + 1 = 0$ gives $K^* = q$, and $C^{\mathrm{LR}}(q) = 2q$. Each
short-run curve touches the long-run line at the output its plant suits,
and lies above it elsewhere: the long-run curve is their lower envelope.
The theorem gives the slope of the envelope without the minimisation:
$\dd C^{\mathrm{LR}}/\dd q = \partial C/\partial q$ at $K = K^*$, that is
$2q/K^* = 2$.

\begin{omfigure}
\begin{tikzpicture}
\begin{axis}[omrecord, grid=none, xlabel={Output $q$}, ylabel={Cost}, xmin=0, xmax=6,
  ymin=0, ymax=14, width=0.62\linewidth, height=0.44\linewidth, samples=60,
  legend pos=north west]
  \addplot[curve, domain=0:4.2] {x^2/1 + 1};
  \addplot[curve, domain=0:5.3] {x^2/2 + 2};
  \addplot[curve, domain=0:6] {x^2/3 + 3};
  \addplot[curve, domain=0:6] {x^2/4 + 4};
  \addplot[supply, domain=0:6] {2*x};
  \addplot[only marks, mark=*, black] coordinates {(1,2) (2,4) (3,6) (4,8)};
\end{axis}
\end{tikzpicture}

{\small Short-run cost curves $q^2/K + K$ for plants $K = 1, 2, 3, 4$
(black), and their lower envelope, the long-run cost $2q$ (red). Plant $K$
touches the envelope at $q = K$.}
\end{omfigure}

The same theorem, applied to the consumer, the cost-minimising firm and the
profit-maximising firm, gives three results used in the next part of the
book: Roy's identity (\cref{ch:b1:consumer-choice}), Shephard's lemma
(\cref{ch:b1:cost-minimisation}) and Hotelling's lemma
(\cref{ch:b1:profit-maximisation}).

\section{Sign predictions and their tests}\label{sec:b1:comparative-statics:signs}

Many models do not pin down the size of a response, only its direction.

\begin{definition}[Sign prediction]\label{def:b1:comparative-statics:sign-prediction}
\index{sign prediction}
A \emph{sign prediction} is a comparative-statics statement that a
derivative is positive, negative or zero, derived from assumptions that
restrict the signs of a model's derivatives without fixing their sizes.
\end{definition}

\Cref{prop:b1:comparative-statics:incidence} is a \omterm{def:b1:comparative-statics:sign-prediction}{sign prediction}: whatever
the slopes, a tax raises the buyers' price by less than the tax and lowers
the quantity. The second derivatives of an objective give another.

\begin{proposition}[Monotone comparative statics]\label{prop:b1:comparative-statics:monotone}
\index{comparative statics!monotone}
Let $x^*(\theta)$ be an interior maximiser of $f(x, \theta)$ with
$f_{xx} < 0$ at the optimum. Then $\dd x^*/\dd\theta$ has the sign of the
cross derivative $f_{x\theta}$.
\end{proposition}

\begin{proof}
The first-order condition is $F(x, \theta) = f_x(x, \theta) = 0$. By
\cref{prop:b1:comparative-statics:implicit}, $\dd x^*/\dd\theta =
-f_{x\theta}/f_{xx}$, and $-1/f_{xx} > 0$.
\end{proof}

A firm choosing output $q$ to maximise $pq - c(q)$ has $f_{qp} = 1 > 0$, so
output rises with the price whatever the cost function, provided the
second-order condition holds. A \omterm{def:b1:comparative-statics:sign-prediction}{sign prediction} is weaker than a number and
easier to refute: one clean case of a tax that raises the price by more
than the tax, in a market that fits the assumptions, contradicts
\cref{prop:b1:comparative-statics:incidence}. That is how the evidence on
pass-through is read.

\begin{model}{A market with a per-unit tax}
\assumes Many buyers and sellers who take the price as given; smooth demand
$D(p)$ and supply $S(p - \tau)$, with $D' < 0 < S'$; buyers and sellers react
to the tax as to any other change in price.
\predicts The price paid by buyers rises by the fraction
$\varepsilon_s/(\varepsilon_s - \varepsilon_d)$ of the tax, between 0 and 1;
the quantity falls by more when demand is more elastic.
\tested American state gasoline and diesel taxes are, on average, fully
passed on to pump prices, which fits a supply to one state that is close to
perfectly elastic; when refineries run above 95 per cent of capacity, only
half of a diesel tax is passed on (Marion and Muehlegger 2011). % ledger: b1:marion2011-pass
The short-run price elasticity of American gasoline demand was $-0.21$ to
$-0.34$ in 1975--1980 and $-0.034$ to $-0.077$ in 2001--2006 (Hughes,
Knittel and Sperling 2008). % ledger: b1:hughes2008-elast
\breaks When sellers have market power, pass-through can exceed one
(\cref{exo:b1:comparative-statics:monopoly}); when a tax is not shown in the
posted price, buyers react to it less than to an equal rise in the price
(the record box below).
\end{model}

\begin{record}[Salience]
For three weeks, a grocery store posted prices that included the sales tax
on 750 of its products, instead of adding the tax at the till (Chetty,
Looney and Kroft 2009). % ledger: b1:chetty2009-salience
The tax was the same; only its display changed. Demand for those products
fell by about 8 per cent relative to comparable products and nearby stores. % ledger: b1:chetty2009-salience
The same study found that rises in alcohol excise taxes, which are included
in posted prices, reduced consumption more than rises in sales taxes, which
are added at the till. The model's assumption that buyers treat a tax like
any other part of the price fails when the tax is out of sight.
\end{record}

\section*{Exercises}\label{sec:b1:comparative-statics:exercises}

\begin{exercise}[$\star$]\label{exo:b1:comparative-statics:ce}
Demand is $q = 50p^{-0.5}$.
\begin{enumerate}
  \item Compute its elasticity at any price.
  \item At $p = 4$, compute $q$; predict the quantity after a $10\%$ rise in price, using the elasticity, and compare with the exact value.
\end{enumerate}
\end{exercise}

\begin{exercise}[$\star$]\label{exo:b1:comparative-statics:read-slope}
In the opening chart the slope of the line is 1.25. % figdata: energy_income_stats
\begin{enumerate}
  \item By what factor does the line say energy use per person differs between two countries, one twice as rich as the other?
  \item India (\usd{7702}, 21.85 million Btu) and the United States (\usd{69121}, 292.7): compute the slope of the segment joining them on the chart. % figdata: energy_income_2019
\end{enumerate}
\end{exercise}

\begin{exercise}[$\star$]\label{exo:b1:comparative-statics:demand-shift}
In the market model $q = a - bp + u^d$, $q = c + dp + u^s$
(\cref{prop:b1:models-and-data:linear}), find the sign of
$\dd q^*/\dd a$ and of $\dd p^*/\dd a$ with \cref{prop:b1:comparative-statics:implicit},
writing the \omterm{def:b1:models-and-data:equilibrium}{equilibrium} as one equation in $p$.
\end{exercise}

\begin{exercise}[$\star$]\label{exo:b1:comparative-statics:spending}
Show that the elasticity of spending $pq$ with respect to the price is
$1 + \varepsilon_d$. If $\varepsilon_d = -0.3$, by how much does spending change
after a $10\%$ rise in price, to first order?
\end{exercise}

\begin{exercise}[$\star\star$]\label{exo:b1:comparative-statics:implicit-tax}
With $D(p) = 100 - 2p$ and $S(p - \tau) = 3(p - \tau)$:
\begin{enumerate}
  \item write $F(p, \tau) = 0$ and compute $F_p$ and $F_\tau$;
  \item compute $\dd p/\dd\tau$, $\dd(p - \tau)/\dd\tau$ and $\dd q/\dd\tau$;
  \item check them against the explicit solution.
\end{enumerate}
\end{exercise}

\begin{exercise}[$\star\star$]\label{exo:b1:comparative-statics:fuel}
The short-run price elasticity of fuel demand is $-0.3$ and the elasticity
of supply to one region is $1.2$.
\begin{enumerate}
  \item Compute the pass-through of a fuel tax.
  \item Predict the rise in the pump price after a tax rise of ten cents a litre, and the share borne by sellers.
  \item Predict how the answer changes if supply becomes far more elastic.
\end{enumerate}
\end{exercise}

\begin{exercise}[$\star\star$]\label{exo:b1:comparative-statics:envelope-simple}
Let $V(a) = \max_x (ax - x^2)$. Find $x^*(a)$ and $V(a)$, and check the
envelope theorem: $V'(a) = \partial f/\partial a$ at $x^*$.
\end{exercise}

\begin{exercise}[$\star\star$]\label{exo:b1:comparative-statics:two-years}
On the 152 countries present in both 1990 and 2019, the slope of the line
is 1.35 in 1990 and 1.25 in 2019. % figdata: energy_income_common
\begin{enumerate}
  \item What does a fall in the slope say about the cross-section of the world?
  \item Why compare the slopes on the same countries?
  \item Does the fall say that a country's energy use now responds less to its income growth? Say what would.
\end{enumerate}
\end{exercise}

\begin{exercise}[$\star\star$]\label{exo:b1:comparative-statics:which-curve}
Is the slope of the opening chart a demand elasticity, a supply
elasticity, or neither? Name the variation the chart uses and two
assumptions that would have to hold for the slope to measure how a given
country's energy demand responds to its income.
\end{exercise}

\begin{exercise}[$\star\star\star$]\label{exo:b1:comparative-statics:envelope-proof}
Prove part (2) of \cref{thm:b1:comparative-statics:envelope}: differentiate
$V(\theta) = f(x^*(\theta), y^*(\theta), \theta)$ and the identity
$g(x^*(\theta), y^*(\theta), \theta) = 0$, and use the \omterm{thm:b1:constrained-optimisation:lagrange}{Lagrange conditions}.
\end{exercise}

\begin{exercise}[$\star\star\star$]\label{exo:b1:comparative-statics:monopoly}
A single seller faces demand $q = Ap^{-\varepsilon}$ with $\varepsilon > 1$ and
has a constant marginal cost $c$.
\begin{enumerate}
  \item Show that its profit-maximising price is $p = c\,\varepsilon/(\varepsilon - 1)$.
  \item A per-unit tax $\tau$ is levied on it. Show that the pass-through $\dd p/\dd\tau$ is $\varepsilon/(\varepsilon - 1) > 1$.
  \item With $\varepsilon = 2$, $c = 10$ and $\tau = 1$, compute the prices before and after.
\end{enumerate}
\end{exercise}

\begin{exercise}[$\star\star\star$]\label{exo:b1:comparative-statics:refute}
\begin{enumerate}
  \item Describe an observation that would refute ``the pass-through of a tax lies between 0 and 1'' in a competitive market, and one that would not refute it although it looks as if it does.
  \item The salience study found that posting the tax lowered demand. Which assumption of the model of the tax does that contradict, and does it change the prediction of pass-through for a tax that is posted?
\end{enumerate}
\end{exercise}

\section{Problem: A Levy on Cement}\label{sec:b1:comparative-statics:problem}

\begin{problem}[{Weekend problem --- a carbon levy on Kestria's cement}]\label{pb:b1:comparative-statics:cement}
Kestria (invented) produces and uses about a million tonnes of cement a
year. Its government plans a carbon levy of $\tau$ kests a tonne, collected
from producers. Buyers pay $p$ a tonne, producers keep $p - \tau$, and $Q$ is in
thousands of tonnes. Demand and supply are
\[
D(p) = 1000\Bigl(\frac{p}{100}\Bigr)^{-0.5}, \qquad
S(p - \tau) = 1000\Bigl(\frac{p - \tau}{100}\Bigr)^{1.2}.
\]

\medskip\noindent\textbf{Part I --- The equilibrium and its derivatives.}
\begin{enumerate}
  \item Show that with no levy the market clears at $p = 100$ kests and $Q = 1000$.
  \item Write the \omterm{def:b1:models-and-data:equilibrium}{equilibrium} condition with the levy as $F(p, \tau) = 0$.
  \item Compute $D'(100)$ and $S'(100)$, then $F_p$ and $F_\tau$ at $(100, 0)$.
  \item Compute $\dd p/\dd\tau$ and $\dd Q/\dd\tau$ at $\tau = 0$, and give their units.
  \item Compute $\dd(p - \tau)/\dd\tau$ and say what it measures.
  \item Show that the straight lines $Q = 1500 - 5p$ and $Q = 12(p - \tau) - 200$ are tangent to $D$ and $S$ at $(100, 1000)$, and solve the linear model for $p$ as a function of $\tau$.
\end{enumerate}

\medskip\noindent\textbf{Part II --- Elasticities and pass-through.}
\begin{enumerate}[resume]
  \item Give the elasticities of demand and of supply, using \cref{def:b1:comparative-statics:constant-elasticity}.
  \item Compute the pass-through with \cref{prop:b1:comparative-statics:pass-through}, to four decimals, and check it against question 4.
  \item What share of the levy do producers bear, to first order?
  \item Name two changes in the market that would push the pass-through towards one.
  \item Compute the elasticity of buyers' spending $pD(p)$ with respect to $p$.
  \item Predict, to first order, the percentage change in buyers' spending on cement after a levy of 10 kests.
\end{enumerate}

\medskip\noindent\textbf{Part III --- What producers lose.}
The producers' total profit, as a function of the price they keep, is
$\Pi(p_s) = \max_Q\, p_s Q - C(Q)$, where $C$ is the industry's cost and
$C'(Q) = 100(Q/1000)^{1/1.2}$.
\begin{enumerate}[resume]
  \item Show that the profit-maximising $Q$ at the price $p_s$ is $S(p_s)$, and that $\Pi'(p_s) = S(p_s)$ by the envelope theorem.
  \item Using question 5, give the first-order fall in the price producers keep after a levy of 10 kests, and the first-order loss of profit.
  \item Compute the first-order loss of the buyers, $Q\,\Delta p$, and the sum of the two first-order losses. Compare it with the levy's revenue at $Q = 1000$.
  \item Will the exact loss of profit be larger or smaller than the first-order one? Explain from the shape of $\Pi$.
\end{enumerate}

\medskip\noindent\textbf{Part IV --- A levy of 10 kests a tonne.}
\begin{enumerate}[resume]
  \item Give the first-order predictions of $p$, $p - \tau$ and $Q$ (to one decimal), using Parts I and II.
  \item The exact \omterm{def:b1:models-and-data:equilibrium}{equilibrium} has $p = 107.16$. Check that it clears the market to three significant figures, and compute $Q$ and $p - \tau$.
  \item Compare the first-order predictions with the exact values. Which is larger, and by how much?
  \item Compute the levy's revenue at the exact \omterm{def:b1:models-and-data:equilibrium}{equilibrium}.
  \item Compute buyers' spending before and after, and its exact percentage change. Compare with question 12.
  \item The exact loss of profit is $\int_{97.16}^{100} S(s)\,\dd s$. Compute it and compare with question 14.
  \item Does the linear model of question 6 give the first-order prediction or the exact one? Why?
  \item State the chapter's result for this levy: the pass-through, in kest cents per kest of levy, to first order and exactly.
\end{enumerate}
\end{problem}
